\section{Discussion}
\label{sec:discussion}

The results raise a few questions that we want to discuss before turning to
the threats to validity.

\paragraph{Is a larger grammar always better?} Not for a wrong-result fuzzer. A
larger grammar reaches more corner cases of SQL semantics, and some of them
break the assumptions an oracle makes: \texttt{HAVING} without \texttt{GROUP
BY}, \texttt{DISTINCT} over \texttt{numeric} values of different scale, or
\texttt{interval} values that are equal but print differently. The price is paid
in oracle precision, and every rewrite needs a soundness check on the tree.

\paragraph{Why keep the tree?} Mutation re-derives a subtree, the oracles
rewrite one typed node, and reduction replaces a subtree with its smallest
derivation. In each case the generator replays the tree and re-binds names and
types, so the result stays valid SQL. On the \texttt{JSON\_TABLE} bug, tree
reduction followed by token reduction gave a 79-token reproducer, whereas token
reduction alone stopped at 116 tokens. On the tree, a clause is one node with a
known span, which lets \tool apply TLP inside a derived table or a
\texttt{LATERAL} subquery, where a text-based rewrite would have to parse the
query again.

\paragraph{Why compare with the last release?} Plan variants cannot catch a
wrong result that every plan shares, such as a constant that is folded before
any plan is chosen (Figure~\ref{fig:motivating}). The last release does not have
the newest planner rules, so it disagrees exactly on the bugs that new code
introduces. It also disagrees on intended changes, which triage has to
recognise.

\paragraph{Can an LLM be trusted?} Not on its word. A rewrite is a claim about
semantics, and \rqref{4} shows that our gate accepted a wrong one. For
fragments the situation is easier, because the model only has to name SQL that
reaches a commit, and the \gate rejects anything that does not parse, is never
derived, fails on the server or misses the changed lines. Repair rounds helped
only rarely: of the 4 accepted commits, 2 passed in the first round,
and 6 of the 10 commits never passed.

\paragraph{Does the approach carry over to other engines?}
\label{sec:beyond}
After the \pg campaigns we wanted to know whether the same design works beyond
\pg, so we generalized \tool to DuckDB, SQLite, Turso and Apache DataFusion
(branch \texttt{generalized} of our repository). The answer is yes, with a small
adapter per engine. Each engine runs in its own worker process, and a dialect
\emph{overlay} of a few dozen lines drops or rewrites the parts of the \pg
grammar that the engine does not accept. With hand-written overlays, 96.3\% of
the generated queries are valid on DuckDB, 97.7\% on SQLite, 93.3\% on Turso
and 79.4\% on DataFusion. An overlay miner can also draft the overlay from the
measured failure rate of each grammar alternative: one mined round lifted DuckDB
from 32.1\% to 44.1\% valid statements without any line written by hand. Instead of the
\pg plan variants, \tool uses optimizer-rule firing, read from rule switches or
from per-rule plan traces, as feedback and as a coverage measure. We also added rows synthesized with Z3 that make TLP
partitions non-empty, and normalizers learned from probes for cross-engine
comparison; between DuckDB and SQLite, these normalizers reduce held-out
disagreements from 28.7\% to 3.6\%. The engines' own test suites serve as seeds,
and a loop triages new failures and drafts a filing package with a regression
test in the engine's own format.
Finally, \emph{rule kernels}, minimal queries that each make one optimizer rule
fire, are composed in pairs. In 15-minute campaigns (medians of three trials per
arm), they raise the number of distinct optimizer-rule pairs reached from 307 to
386 on DuckDB, from 241 to 276 on DataFusion, and from 86 to 182 on SQLite.
Every kernel trial beat every plain one ($\hat{A}_{12}=1.00$), although three
trials per arm cannot give $p<0.1$.

These campaigns found fourteen bugs. Nine of them were already fixed on the
engine's development branch when we checked. For the other five we filed an
issue and a pull request with a fix and a regression test: Turso \#11594 / PR
\#11595 (\texttt{percentile\_disc} with an expression argument) and \#11796 / PR
\#11797 (an aggregate in \texttt{IN (SELECT ...)} next to a window function,
which gives a spurious error or a panic); DuckDB \#28551 / PR \#28552
(\texttt{min}/\texttt{max} over a nested value computed by a subquery, a 2.0.0
regression), which a maintainer approved for merging and for the 2.0 release
line; DataFusion \#27235 / PR \#27236 (\texttt{UPDATE} and \texttt{DELETE}
statements that repeat an expression fail on an internal common-subexpression
column); and a DataFusion internal error for a scalar subquery inside
\texttt{VALUES}. A single command (\texttt{scripts/reproduce\_generalized.sh})
reproduces all five on the released versions that the campaigns ran on.
